% Folder 146, batch 07 — archivists' pages 121–133.
% Pass: Opus 5.5 (claude-opus-5-5), 2026-10-02 — first pass, unchecked against the pages by a human.
%
% Pages 131 and 133 are unrelated typed versos (USTL thesis-defence notices,
% November 1981) carrying nothing in his hand, and are skipped. Page 125 is a
% cover sheet in his hand (« Graphes MD »); it gets no \page and its words
% become the section heading.
\documentclass[11pt,a4paper]{article}
\input{../preamble/grothendieck}

\folder{146}
\batch{7}
\pages{121}{133}
\foldertitle{Groupoïdes de Teichmüller et les S0 Tg,!ν, S0 TG : notes manuscrites (s.d.).}
\dating{[à partir de 1978]}
\watermark{Édition de démonstration}

\begin{document}

\section{Cas exceptionnels : torseurs tangents et ensembles $\mathcal{E}_{G,Q}$}

\page{121}
\note{la page commence au milieu d'un argument venu de la page précédente, hors de ce lot : il s'agit des configurations exceptionnelles d'un ensemble $J$ de points sur une droite projective $Y$.}

\uncertain{type} $2,2$ de $J$. Ces points correspondent \uncertain{aux} droites projectives $J$ \ill{} uniquement quand \ill{} des automorphismes \add{d'ordre $>4$} \struck{non triviaux}, il y en a de deux sortes :

\marginal{\uncertain{ils sont} \ill{} d'ordre \ill{} \uncertain{multiples} \ill{} $4$}

a) Pour une structure « carré » sur $J$, correspondant à une partition de type $(1,2)$\note{lecture du type incertaine}, on a un plongement canonique de $Y$ dans un \uncertain{espace projectif}, de façon que les \add{automorphismes} \struck{\uncertain{groupe des}} \struck{\uncertain{automorphismes du}} carré soient induits par des automorphismes projectifs. Si $Q$ est une structure carré sur $J$, on désigne par $Y_Q$ la droite projective correspondante. \struck{\ill{}}

\drawing{une courbe fermée (ovale) portant quatre points marqués, en regard des lignes précédentes.}

Le groupe des rotations du carré opérant de façon simplement transitive sur l'ensemble $J$, de sorte que les $T_i$ s'identifient entre eux. Si $\omega_Q$ est l'une des deux orientations du carré, \struck{qui est \ill{} \ill{} $1$-$1$ \ill{} \ill{} l'une des deux orientations de l'ens.\ $J$,} on a un isomorphisme canonique
\[
T_i = \mathbb{U} \times_{\{\pm1\}} \omega_Q
\]
en choisissant comme \add{en $i$} origine \ill{} : une orientation \ill{} la direction tangente

\page{122}
\ill{} équivalence, dans le sens de cette orientation.

On aura alors
\[
\text{\struck{$T^{Y}_{\hat{A}_2}$}} = \mathbb{U}^{A}
\]
\[
T^{Y}_{J_{\mathrm{I}}} = \mathbb{U}^{J_{\mathrm{I}}} \wedge_{\{\pm1\}} \omega_Q
\]
\marginal{NB $\{\pm1\}$ s'envoie diagonalement dans $\mathbb{U}^{J}$,}
\[
T^{Y}_{J_{\mathrm{II}}} \simeq \mathbb{U}^{J_2}
\]
car $\bigwedge_{i\in J_{\mathrm{II}}} T_i \simeq \mathbb{U}$, puisque \add{\uncertain{somme}} les $T_i$ sont isomorphes \uncertain{canon.}\ à $\mathbb{U}\wedge_{\pm1}\omega_Q$, dans le \uncertain{produit contracté} \ill{} \uncertain{d'autre} \ill{} \ill{} $\mathbb{U}$.
\[
T^{Y}_{J_{\mathrm{III}}} \simeq \prod_{i\in J_{\mathrm{III}}} \mathbb{U}\wedge_{\{\pm1\}}
\underbrace{(\omega_Q \wedge \tilde{J}_{\mathrm{III}\,i})}_{\mathcal{E}_i}
\]
Donc, si $J_{\mathrm{I}}\neq\emptyset$, on introduit l'ensemble
\[
\mathcal{E}_{G,Q} = \mathcal{E}_{G,Y_Q}
= \underbrace{\mathcal{E}_{\hat{A}_1}}_{\prod_{a\in\hat{A}_1}\mathcal{E}_a}
\times \omega_Q \times
\underbrace{\mathcal{E}_{J_{\mathrm{III}}}}_{\prod_{i\in J_{\mathrm{III}}}\omega_Q\wedge\tilde{J}_{\mathrm{III}\,i}}
\]
qui est un torseur sous
\[
\gamma_{G_{\uncertain{1}}} = \{\pm1\}^{\hat{A}_1}\times\{\pm1\}\times\{\pm1\}^{J_{\mathrm{III}}}
\]
\note{l'indice de $\gamma$ est peu lisible ; la suite écrit $\gamma_G$.}
qui \uncertain{s'envoie} dans $\mathbb{U}^{\hat{A}}$
\[
\gamma_G \longrightarrow \mathbb{U}^{\hat{A}} \simeq \mathbb{U}^{\hat{A}_1}\times\mathbb{U}^{J_{\mathrm{I}}}
\times\mathbb{U}^{J_{\mathrm{II}}/\sigma}\times\mathbb{U}^{J_{\mathrm{III}}}
\]
et (flèche directe, tracée sous la précédente)
\[
\gamma_G \longrightarrow \mathbb{U}^{\hat{A}_1}\times\check{\mathbb{U}}^{J_{\mathrm{I}}}\times\mathbb{U}^{J_{\mathrm{III}}}
\]
\add{\uncertain{comme}} se factorisant par ce produit triple, et en

\page{123}
prenant le produit cartésien de
\[
\{\pm1\}^{A_1} \hookrightarrow \mathbb{U}^{A_1}
\qquad\text{homom.\ évident déduit de } \{\pm1\}\hookrightarrow\mathbb{U}
\]
\[
\{\pm1\} \longrightarrow \mathbb{U}^{J_{\mathrm{I}}},
\qquad
\{\pm1\}\xrightarrow{\ \mathrm{diag}\ }\{\pm1\}^{J_{\mathrm{I}}}
\ \text{\uncertain{$\Sigma$ t.\ d.\ car}}
\]
\[
\{\pm1\}^{J_{\mathrm{III}}} \hookrightarrow \mathbb{U}^{J_{\mathrm{III}}}
\qquad\text{homom.\ évident.}
\]
Ceci posé, on aura
\[
SP^{Y}_{G} \simeq \mathbb{U}^{\hat{A}} \wedge_{\gamma_G}
\mathcal{E}_{G,Q} \xrightarrow{\ \gamma_G\text{-tors.}\ }
\]
\note{le but de la flèche n'est pas écrit ; sous la flèche : « $\gamma_G$-tors. ».}

Dans le cas où $J_{\mathrm{I}}=\emptyset$, on \add{peut} poser \ill{} \uncertain{pareil}
\[
\mathcal{E}_{G,Q} = \mathcal{E}_{G,Y_Q} = \mathcal{E}_{\hat{A}_1}\times\mathcal{E}_{J_{\mathrm{III}}}
\]
\uncertain{puisque} l'\uncertain{économie} des \uncertain{facteurs} $\omega_Q$, \ill{}
\[
\gamma_G = \{\pm1\}^{\hat{A}_1}\times\{\pm1\}^{J_{\mathrm{III}}} = \{\pm1\}^{\hat{A}_1\amalg J_{\mathrm{III}}}
\]
et \ill{} \ill{} l'expression ci-dessus pour $SP^{Y}_{G}$.

Considérons alors
\[
S_0T^{Y_Q}_{G} = S_0T^{Q}_{G}
\]
son groupoïde plein de $\pi_1(SP^{Y}_{G})$ \struck{$\Pi$}\note{le $\pi_1$ est récrit en surcharge sur un autre symbole} basé en les éléments de $\mathcal{E}_{G,Q}$. \uncertain{Les} formules explicites de ce groupoïde \add{basé en $\mathcal{E}_{G,Q}$} \struck{\uncertain{fondamental}} \uncertain{par} générateurs et relations --- j'y reviendrai.

\page{124}
Je vais d'abord passer en revue les cas de $Y$ exceptionnels :

Cas où il existe $i_0\in J$ tel que les \uncertain{autres} \ill{} de $(Y,J)$ fixant $i_0$ forment un groupe d'ordre $3$ \struck{\uncertain{Cas $Y$} \ill{} \ill{} \uncertain{le domaine} \ill{} \ill{} \uncertain{triangulaire} \ill{} [$i_0\in J$] et de \ill{} \ill{} \uncertain{circulaire}} [Je reviendrai \ill{} \ill{} dire que le groupe \uncertain{alterné} \ill{} \ill{}] \ill{}

\drawing{à gauche, deux figures : un cercle (la sphère) traversé d'arcs, avec un point $i_0$ en haut et trois points marqués ; plus bas, un triangle de sommets $a$, $b$, $c$ avec un point au-dessus du sommet supérieur.}

$(Y,J)$ est le sous-groupe $\mathfrak{S}^{+}_{J}$ de \ill{} $\mathfrak{S}_J$ \ill{} \ill{} \ill{} \ill{} que la configuration \ill{} \ill{} \uncertain{de la sphère} (orientée) \ill{} les 2 possibilités \ill{} \uncertain{orientations} \ill{} \ill{} \ill{} l'un des \ill{} \uncertain{seul} \ill{} de $J$, \ill{} \ill{}. Le groupe d'isotropie est isomorphe \ill{} d'une paire \ill{} \ill{} circulaires sur l'un des \ill{} des permutations \uncertain{circulaires} des trois autres points, qui forment un torseur sous $\mu_3=\mu_3(\mathbb{C})\simeq\mathbb{Z}/3\mathbb{Z}$, soit $\eta_i$ ($i\in I$). On a alors
\[
J_i \simeq \mathbb{U}\times_{\mu_3}\eta_i
\]
où $\mu_3\hookrightarrow\mathbb{U}$ est l'inclusion canonique.

On \ill{}

\note{la page s'arrête sur « On » ; la suite manque dans ce lot.}

\section{Graphes MD}
\note{titre de sa main, sur un feuillet par ailleurs blanc (p.~125) ; la mention « (6 p) » au crayon n'est pas de lui.}

\page{126}
\drawing{tableau de graphes, une ligne par type $(g,\nu)$, avec à droite de chaque ligne un nombre qui est vraisemblablement la dimension $3g-3+\nu$. Dans chaque ligne, les graphes MD du type sont dessinés (sommets marqués de leur genre quand il est non nul, arêtes libres en flèches sortantes), reliés par des flèches de spécialisation dirigées vers les graphes plus simples, certaines affectées d'une multiplicité ($2$, $3$).
$(0,3)$ : un sommet à trois arêtes libres ; nombre $0$.
$(0,4)$ : le sommet à quatre arêtes libres $\leftarrow$ deux sommets trivalents joints par une arête ; nombre $1$.
$(0,5)$ : sommet à cinq arêtes libres $\leftarrow$ deux sommets $\xleftarrow{2}$ trois sommets en chaîne ; nombre $2$.
$(0,6)$ : sommet à six arêtes libres, puis graphes à deux, trois et quatre sommets, dont une étoile à trois branches, avec flèches marquées $2$ et $3$ ; nombre $3$.
$(1,1)$ : un sommet de genre $1$ à une arête libre $\leftarrow$ une boucle à une arête libre ; nombre $1$.
$(1,2)$ : sommet de genre $1$ à deux arêtes libres, puis boucle à deux arêtes libres, sommet de genre $1$ relié à un sommet trivalent, deux sommets joints par une double arête, etc., flèche marquée $2$ ; nombre $2$.
$(1,3)$ : diagramme analogue, plus touffu, aboutissant à un triangle à trois arêtes libres, flèches marquées $2$ et $3$ ; nombre $3$.
$(2,0)$ : sommet de genre $2$ ; sommet de genre $1$ à une boucle ; deux sommets de genre $1$ joints par une arête ; sommet de genre $1$ relié à une boucle ; huit (deux boucles en un sommet) ; deux sommets joints par trois arêtes ; flèches marquées $2$, $3$ et une flèche marquée d'un signe peu lisible ; nombre $3$.
$(2,1)$ : premiers graphes seulement, puis « ……… » ; nombre $4$.}
\note{un graphe du type $(1,3)$ et un du type $(2,0)$ sont noircis (biffés) sur la feuille.}

\page{127}
\emph{Graphes MD}

Un \emph{graphe pondéré} est donné par
\[
(1)\qquad (S,\hat{\vec{A}},\sigma,\bar{\sigma},g)
\]
où $S$, $\hat{\vec{A}}$ ens.\ finis, $\sigma:\hat{\vec{A}}\to S$ application (« origine »), $\bar{\sigma}:\hat{\vec{A}}\to\hat{\vec{A}}$ une involution, $g:S\to\mathbb{N}$ une application :
\[
(2)\qquad
\underset{\text{arêtes}}{\hat{\vec{A}}}\xrightarrow{\ \sigma\ }\underset{\text{sommets}}{S}\xrightarrow{\ g\ }\mathbb{N}
\qquad \bar{\sigma}^2=\mathrm{id}_{\hat{\vec{A}}}
\]
\note{sous la flèche $\sigma$ : « appl.\ origine » ; sous $g$ : « application “genre” » ; l'involution $\bar\sigma$ est dessinée comme une boucle sur $\hat{\vec{A}}$, avec la légende « renversement \uncertain{d'orientation} ».}

On pose
\[
\begin{cases}
\hat{\vec{A}}^{\bar\sigma} = I \hookrightarrow \hat{A}\\[2pt]
\vec{A} = \hat{\vec{A}}\smallsetminus\hat{\vec{A}}^{\bar\sigma}\\[2pt]
A = \vec{A}/\bar\sigma \simeq \hat{A}\smallsetminus I
\end{cases}
\]
Donc
\[
(3)\qquad \hat{\vec{A}}/\bar\sigma = \hat{A} \qquad (\text{arêtes})
\]
\add{\uncertain{non orientées}} arêtes \emph{ordinaires} (\uncertain{i.e.}\ les arêtes correspondant à : deux éléments), les autres sont dites \emph{\uncertain{bordantes}} \struck{\uncertain{relatives}} \add{\uncertain{ou}} arêtes libres. Représentations géographiques
\add{NB On ne marque pas $g_i$ si $g_i=0$}

\drawing{un graphe à trois sommets $g_1$, $g_2$, $g_3$ : $g_1$ porte trois arêtes libres et deux arêtes vers $g_2$, $g_2$ porte une boucle et une arête libre et est relié à $g_3$, qui porte deux arêtes libres.}

\marginal{graphes \emph{stables} : \ill{} \uncertain{complètement} \ill{} \uncertain{ordinaires}, \uncertain{rendant} \ill{} \ill{} qui satisfont (4)}

Condition de stabilité MD :
\[
(4)\qquad \forall s\in S,\quad 2g_s+\hat{\nu}_s\geq 3
\quad\text{i.e.}\quad
\begin{cases}
g_s=0 \Rightarrow \hat{\nu}_s\geq 3\\
g_s=1 \Rightarrow \hat{\nu}_s\geq 1
\end{cases}
\]
\note{l'inégalité est écrite $2g_s+\hat\nu_s\geq3$, sans le cas $g_s\geq 2$, qui ne demande rien.}
où $\hat{\nu}_s$ est le nb d'arcs \struck{qui} d'origine $s$ (\emph{ordre total} du sommet $s$) et \ill{} par $\nu_s$ désignons le nb d'arcs \add{$\hat{\vec{A}}^{\bar\sigma}$ i.e.\ $I=\hat{\vec{A}}^{\bar\sigma}$} \ill{} \uncertain{dans} \uncertain{card} $I$
\[
(5)\qquad \hat{\nu}_s = \nu_s + \vec{\nu}_s ,
\qquad \text{\struck{\ill{}}}
\]
\note{sous les termes : « ordre total de $s$ », « ordre \struck{\ill{}} \uncertain{[libre]} de $s$ », « ordre complémentaire de $s$ \uncertain{correspondant aux} “pts marqués” \uncertain{sur une courbe} \ill{} » --- et « \uncertain{marqué} ». L'affectation de chaque légende à $\hat\nu_s$, $\nu_s$, $\vec\nu_s$ est incertaine.}

\marginal{ordinaires \ill{} \ill{} \ill{} \ill{} $\vec{A}/\bar\sigma\to$ \ill{} $\pi_0(G)\to\mathbb{N}$ \ill{} \uncertain{genre} \ill{} \ill{}}
\marginal{\ill{} le graphe \ill{} $(g,\nu)$ \ill{} \uncertain{Weil} \ill{} \ill{}}
\note{deux notes marginales en oblique, cernées d'un trait, à gauche du bas de la page ; je n'en lis que des fragments.}

On pose
\[
(6)\qquad \nu=\nu(G) = \sum_s \vec{\nu}_s = \operatorname{card}(\vec{A}^{\bar\sigma})
\quad(\text{\emph{ordre marqué} \struck{total} du graphe})
\]
\[
(7)\qquad c=c(G)=\tfrac12\sum_s\nu_s
\quad\text{nb des arêtes ordinaires : \emph{codim.\ modulaire} du graphe}
\]
\[
(8)\qquad g=\sum_{s\in S}g_s+\underbrace{h^1(G)}_{\text{connexité de }G}
\quad\text{\emph{genre virtuel} du graphe}
\]
(correspond à $\dim H^1(X,\mathcal{O}_X)$)
\note{le premier membre de (7) porte un chiffre biffé avant $\tfrac12$ ; la confusion entre $\nu_s$ et $\vec\nu_s$ dans (5)--(7) est celle de la page.}

\page{128}
Si le graphe est MD, alors $(g,\nu)$ est \uncertain{admissible}, et on pose
\[
(9)\qquad \delta(G)=3g-3+\nu
\qquad(\text{\emph{dimension modulaire}, \add{\emph{\uncertain{composante}}}})
\]
et on aura
\[
(10)\qquad \delta(G)+c(G)=\delta(\mathbf{G})
\]
\note{le second membre est écrit avec un $G$ d'une autre forme, sans doute le type $(g,\nu)$ du graphe.}

\emph{Opérations} (Soit $G,G'$ deux graphes MD, on appelle \emph{homomorphismes} \add{de} \emph{\uncertain{spécialisation}} $u$ de $G$ dans $G'$ un couple d'applications
\[
S\xrightarrow{\ u_s\ }S',
\qquad
\hat{\vec{A}}\xrightarrow{\ u_a\ }\hat{\vec{A}}{}'\amalg S'
\]
satisfaisant les conditions suivantes :
\[
\vec{B}=u_a^{-1}(\hat{\vec{A}}{}'),\qquad
\vec{C}=\hat{\vec{A}}\smallsetminus\vec{B}=u_a^{-1}(S')
\]

a) $\vec{B}\supset I$ (ens.\ des arêtes libres), $\vec{B}$ stable par $\bar\sigma$ ;

b) L'application $\vec{B}\to\hat{\vec{A}}{}'$ induite par $u_a$ commute aux involutions et le couple $(u_s,u_a|\vec{B})$ commute aux applications origine. (Donc $u_a$ induit $I\to I'$.)

c) \struck{\ill{}} si $c\in\vec{C}$, alors $u_a(c)=u_s(\sigma(c))$ \add{$=u_a(\bar\sigma c)$}, où $\sigma(c)$ est l'origine (donc $u_a|C$ est connu quand on connaît $u_s$). [Donc $u_s$ \uncertain{est} \struck{\uncertain{déf.\ par} $u_s$} \uncertain{impose} \ill{} deux \struck{\uncertain{couples}} sommets $s,s'$ tels que $\exists c\in\vec{C}$, $s=\sigma(c)$, $s'=\sigma(\bar\sigma c)$ --- donc $u_s$ se factorise en
\[
S\to S/R\to S'
\]
où $R$ est la relation d'équivalence dans $S$ engendrée par la relation précédente.]

d) $S/R\to S'$ et $\vec{B}\to\hat{\vec{A}}{}'$ induit par $u$ est une bijection, est une bijection :

e) Pour tout $s'\in S'$, soit $S(s')=u_s^{-1}(\{s'\})$, $\vec{C}(s')=$ ens.\ des $\alpha\in\vec{C}$ tels que $\sigma(\alpha)$ et $\sigma(\bar\sigma\alpha)\in S(s')$ (i.e.\ $u_a(\alpha)=s'$) \struck{$u_a(\alpha)=s'$ ou $u_a(\alpha)\in B$ et $\sigma(u_a(\alpha))=\ldots$}. Soit le sous-graphe \add{pondéré} $G_0(s')$ \struck{\ill{}} de $G$ défini par $S(s')$ et $\hat{\vec{A}}(s')$, considérons son genre virtuel
\[
g(G_0(s'))=\sum_{s\in S(s')}g_s+h^1(G_0(s'))
\]
On a alors
\[
g_{s'}=g(G_0(s'))
\]

\marginal{i.e.\ $(S,\vec{B})$ \uncertain{défini par} $(S,\hat{\vec{A}},\sigma,\bar\sigma)$ \uncertain{considéré comme} graphe \uncertain{ayant} \ill{} \ill{} $u_s$ \ill{}}
\marginal{d) \uncertain{donc} $S'$ et $\hat{\vec{A}}{}'$ se \uncertain{récupèrent} : $B\subset\hat{\vec{A}}$ \uncertain{bijection} \ill{} \ill{} \uncertain{compris}…}
\marginal{NB d) \uncertain{signifie} \ill{} $G(s')$ \uncertain{connexe}}
\note{la « condition d) » est écrite avec un « et » puis répète « est une bijection » ; le « $S(s')$ » de e) est récrit en surcharge.}

\page{129}
\add{\emph{NB}} Le \add{MD-}graphe \struck{\ill{}} $G'$ se déduit, à isom.\ unique près, de $G$ et de $C$ ($\subset A=A(G)$ \add{$=\vec{C}/\bar\sigma$}, ens.\ des arêtes liées). Soit $G_{0C}$ le sous-graphe de $G$ défini par
\[
S(G_{0C})=S,\qquad
\vec{A}(G_{0C})=\vec{C}\quad(\text{image inv.\ de } C \text{ dans } \vec{A})
\]
pondéré par $g : S=S(G_{0C})\to\mathbb{N}$ provenant de $G$.

On pose
\[
\left\{
\begin{aligned}
&S'=\pi_0(G_{0C})\\
&\hat{\vec{A}}{}'=\hat{\vec{A}}\smallsetminus\vec{C}\\
&\sigma':\hat{\vec{A}}{}'\to S' \text{ composé } \hat{\vec{A}}{}'\hookrightarrow\hat{\vec{A}}\xrightarrow{\ \sigma\ }S=S(G_{0C})\xrightarrow{\ \mathrm{can}\ }S'\\
&\bar\sigma' \text{ induit par } \bar\sigma\\
&g_{s'} = \text{genre virtuel de la comp.\ connexe } G_{0C,s'} \text{ du graphe pondéré } G_{0C}
\end{aligned}
\right.
\]
\note{après $\hat{\vec{A}}\smallsetminus\vec{C}$ un « $=\vec{B}$ » est biffé ; au-dessus de $\hat{\vec{A}}{}'\hookrightarrow\hat{\vec{A}}$ : « où $\vec{C}$ est l'\uncertain{image inverse} » ; sous $S'$ final : « $=\pi_0(G_{0C})$ ».}

\emph{NB} On aura
\[
I\simeq I',\qquad \pi_0(G)\xrightarrow{\ \sim\ }\pi_0(G')
\]
et les applications $\gamma:\pi_0(G)\to\mathbb{N}$ et $\gamma':\pi_0(G')\to\mathbb{N}$ (genre virtuel d'une composante) se correspondent.

\page{130}
1) \struck{Sectionnage} \add{Décampage} (d'arêtes) et recollage (de \uncertain{bouts} d'arêtes libres)

$\Leftrightarrow$ normalisation partielle et recollage (par identification)

Si $G$ et $G'$ se correspondent (\add{$G'$} déduit de $G$ par décampage)
\[
M_G\simeq M_{G'},\qquad \hat{M}_G\simeq\hat{M}_{G'},\qquad
\mathfrak{X}_{G'}\to\mathfrak{X}_G,\quad \hat{\mathfrak{X}}_{G'}\to\hat{\mathfrak{X}}_G
\]
\emph{NB} On a $S=S'$ et $g_s=g'_s$, $\hat{\nu}_s=\hat{\nu}'_s$ --- mais \struck{\ill{}} $\nu_s$, $\vec{\nu}_s$ changent.

Cette opération sert surtout pour exprimer \add{\ill{}} $M_G$ \uncertain{ou} $\hat{M}_G$ comme un produit $\prod_{s\in S}M_{g_s,\hat{\vec{A}}_s}$ \uncertain{resp.}\ $\prod_s\hat{M}_{g_s,\hat{A}_s}$.

2) \add{\uncertain{Homomorphisme}} Hom.\ ordinaire de \add{MD-}graphes \struck{pondérés} (hom.\ de graphes, compatible aux pondérations) $G\xrightarrow{u}G'$ (comme \uncertain{ci-dessus})

\begin{tikzcd}[column sep=small, row sep=small]
S \arrow[rr] \arrow[dr, "g"'] & & S' \arrow[dl, "g'"] \\
& \mathbb{N} &
\end{tikzcd}

\note{le triangle est marqué « commutatif ».}

On a alors
\[
M_G\longrightarrow M_{G'},\qquad
\text{\struck{$\hat{M}_G\to\hat{M}_{G'}$}}
\]
\[
\mathfrak{X}_G\times_{M_{G'}}M_{G'}\longrightarrow\mathfrak{X}_G,
\qquad
\hat{\mathfrak{X}}_G\times_{\hat{M}_{G'}}\hat{M}_G\to\hat{\mathfrak{X}}_G
\]
\note{la seconde formule de chaque ligne est barrée d'une grande croix, avec la note « ne \uncertain{fait} pas \uncertain{sens}, \uncertain{voir} \ill{} \ill{} » ; il écrit au-dessus « $\hat{M}_G\to\hat{\mathfrak{X}}_G$ ». Le texte qui suit est « (car en effet d'ailleurs $M_G\to M_{G'}$, \uncertain{se réduit} \ill{}) ».}

Cas particulier \struck{Surtout \uncertain{intéressant}} intéressant : les homom.\ de transition entre les $M_{g,\nu}$ (qui \uncertain{déjà} \ill{} \ill{} compactifiés !)

\marginal{Mais l'application $\hat{\vec{A}}_{s'}\to\hat{\vec{A}}_s$ ($s=u(s')$) doit être injective --- ou disons \emph{immersion}}

Hom.\ d'induction de MD-graphes $G'$, $G$, $u:G'\to G$
\[
S'\xrightarrow{\ u_0\ }S,\qquad
\hat{\vec{A}}{}'=I'\amalg\vec{A}'\xrightarrow{\ u_1\ }\hat{\vec{A}}=I\amalg\vec{A}
\]
satisfaisant

a) compatibilité aux applications origine

\begin{tikzcd}
\hat{\vec{A}}{}' \arrow[r, "u_1"] \arrow[d, "\sigma'"'] & \hat{\vec{A}} \arrow[d, "\sigma"] \\
S' \arrow[r, "u_0"'] & S
\end{tikzcd}

commutatif (mais pas compatible aux applications extrémité, ni aux $\bar\sigma$) ;

b) $u_1^{-1}(I)\subset I'$ \add{\uncertain{i.e.}} i.e.\ $u_1^{-1}(\vec{A})\supset\vec{A}'$ i.e.\ $u_1(\vec{A}')\subset\vec{A}$, et on veut que sur $\vec{A}'$, $u_1$ commute à $\bar\sigma$, $\bar\sigma'$
\struck{c) \ill{} $I'_1=I'\smallsetminus I'_0=I'\cap u_1^{-1}(\vec{A})$, \ill{} \ill{} que $u_1(I'_1)\subset\vec{A}$ soit stable par $\bar\sigma$, et $u_1|I'_1$ injectif}

\marginal{c) $u_1:\hat{\vec{A}}_{s'}\to\hat{\vec{A}}_s$ ($s=u_0(s')$) \emph{injectif}}
\note{la marge c) remplace la condition c) biffée en bas de page, vers laquelle pointe une flèche.}

\page{132}
\note{la p.~131 est un verso administratif sans rapport ; la condition d) qui suit continue la liste a)--c) de la p.~130.}

d) On a commutativité dans

\begin{tikzcd}[column sep=small, row sep=small]
S' \arrow[rr, "u_0"] \arrow[dr, "g'"'] & & S \arrow[dl, "g"] \\
& \mathbb{N} &
\end{tikzcd}


Soit $X$ une courbe de type $G$, \add{sur un corps $k$}, avec normalisée
\[
\tilde{X}=\coprod_{s\in S}\tilde{X}_s \longleftrightarrow \hat{\vec{A}}_k=\coprod_s(\hat{\vec{A}}_s)_k
\]
$X$ déduit de $\tilde{X}$ par identification à l'aide de $\bar\sigma$ ; on en déduit une courbe $X'$ de type $G'$, \uncertain{posant}
\[
\tilde{X}'=\coprod_{s'\in S'}\tilde{X}'_{s'} \longrightarrow \hat{\vec{A}}{}'_k=\coprod_{s'\in S'}(\hat{\vec{A}}{}'_{s'})_k,
\qquad
\tilde{X}'_{s'}\overset{\mathrm{def}}{=}\tilde{X}_{u_0(s')}
\]
défini via
\[
(\hat{\vec{A}}{}'_{s'})_k\overset{\mathrm{inj}}{\hookrightarrow}(\hat{\vec{A}}_s)_k\hookrightarrow\tilde{X}_s=\tilde{X}'_{s'}
\]
et on construit $X'$ par identification à partir de $\tilde{X}'$ via $\bar\sigma'$, \uncertain{opérant} \add{librement} sur $\vec{A}'_k$. On a évidemment
\[
X'\longrightarrow X \qquad \text{\emph{morphisme fini net}}
\]
appliquant $I'$ dans $I\amalg\operatorname{Sing}X$, et qui sur les \struck{\uncertain{composantes irréd.}} \add{\uncertain{des normalisées}} \struck{\ill{} des normalisées \ill{}} \uncertain{iso}. Inversement \ill{} \ill{} un tel morphisme

---…

\end{document}
